% Pursuing Stacks (cotes 134-2 to 134-8): archive cote 134-2, site page 51 (PDF page 52).
% appendix (letters to Breen). Draft compared: ps1.tex:1254-1278.
% AI-transcribed by Claude Opus 5.5 (claude-opus-5-5, adaptive thinking, effort high), Message Batches API, batch ps1 (msgbatch_01Rv5hCaNTtUt4JHx1yqoHPL), 2026-11-12; not yet checked by eye.
\scanpage{134-2/51}
\mspage{\hand{36}} \mspage{4}

vrai pour le \hand{$\Pi$}$_2$(T). Ainsi, si T est un espace topologique réunion de
deux fermés $T_1$ et $T_2$, il n'est sans doute plus vrai que la donnée d'un 1-champs
localement constant sur T ``équivaut` à\hand{''} \unsure{la} donnée d'un 1-champ localement cons-%
tant \hand{$\mathcal{F}$}$_i$ sur $T_i$ (i = 1,2) et d'une équivalence entre les restrictions
\margin{[bracket from this line to the next]}de \hand{$\mathcal{F}$}$_1$ et \hand{$\mathcal{F}$}$_2$ à $T_1$ \hand{$\cap$} $T_2$ (alors que l'énoncé analogue en termes de
o-champs, i.e. de revêtements, est évidemment correct).

\margin{espace} \margin{[bracket and arrow downward]}\hand{(App.7)} L'énoncé (B) plus haut rend clair comment expliciter la cohomologie
d'un n-groupoide \hand{$\mathcal{C}$}$_n$. Si \hand{$\mathcal{C}$}$_n$ = \hand{$\Pi$}$_n$(X), et si \hand{$\mathcal{F}$} est un (n-1)-champ\xout{\unsure{s}}
localement constant sur X, $e^X_{n-1}$ est le (n-1)-champ ``final'', on a une
(n-1) équivalence de (n-1)- catégories
\[
\hand{$\Gamma$}_X(\hand{$\mathcal{F}$}) = \hand{$\mathcal{F}$}(X) \hand{$\simeq$} \underline{\mathrm{Hom}}\ (e^X_{n-1}, \hand{$\mathcal{F}$})
\]
qui montre que le foncteur \hand{$\Gamma$}$_X$ ``intégration sur X'' sur les (n-1)-champs
localement constants, qui inclut la cohomologie (non commutative) localement
constante de X en dim \hand{$\leq$} n-1, s'interprète en termes de ``(n-1)- systèmes locaux''
sur le groupoide fondamental comme un \uline{Hom} ( $e_{n-1}^{(\hand{$\mathcal{C}_n$})}$ , \hand{$\mathcal{F}$}) où mainte-%
nant \hand{$\mathcal{F}$} est interprété comme un n-foncteur
\[
\hand{$\mathcal{C}$}_n \xrightarrow{\hand{$\mathcal{F}$}} ((n-1)\text{-Cat})
\]
et $e_{n-1}^{(\hand{$\mathcal{C}_n$})}$ est le n-foncteur constant sur \hand{$\mathcal{C}$}$_n$, de valeur la (n-1)-catégorie
finale.

Pour interpréter ceci en notation cohomologique, il faut que j'ajoute,
comme ``remords'' à la lettre précédente, l'interprétation explicite de la
cohomologie non commutative sur un topos X, en termes d'intégration de
n-champs sur X. Soit \hand{$\mathcal{F}$} un n-champ de Picard strict sur X, il est donc
défini par un complexe de cochaines \hand{$\mathcal{L}^{\cdot}$} sur X
\[
\hand{$0 \to \mathcal{L}^0 \to \mathcal{L}^1 \to \mathcal{L}^2 \to \dots \to \mathcal{L}^n \to 0$}
\]
concentré en degrés o \hand{$\leq$} i \hand{$\leq$} n (défini à isomorphisme\hand{/}unique près dans la
catégorie dérivée de Ab(X)\hand{)}. Ceci dit, les $H^i$ (X, \hand{$\mathcal{L}^{\cdot}$} ) (hypercohomologie)
\uline{pour 0 \hand{$\leq$} i \hand{$\leq$} n} s'interprêtent comme $H^i$(X,\hand{$\mathcal{L}^{\cdot}$}) = \hand{$\pi$}$_{n-i}$ \hand{$\Gamma$}$_X$(\hand{$\mathcal{F}$}).
Si on s'intéresse à tous les $H^i$ (pas seulement pour i \hand{$\leq$} n) on doit, pour
tout N \hand{$\geq$} n, regarder \hand{$\mathcal{L}^{\cdot}$} comme un complexe concentré en degrés 0 \hand{$\leq$} i \hand{$\leq$} N

\margin{51}
